\section{Matematica}

\subsection{Identidades}

\subsubsection{Series}

$$\sum_{i=1}^{n} i^{2} = \frac{n(n+1)(2n+1)}{6}  \qquad  \sum_{i=1}^{n} i^{3} = \frac{n^{2}(n+1)^{2}}{4} = \left(\sum_{i=1}^n i\right)^2$$

$$ g_k(n) = \sum_{i=1}^n i^k = \frac{1}{k+1} \left( n^{k+1} + \sum_{j=1}^k \binom{k+1}{j+1} (-1)^{j+1} g_{k-j}(n) \right) $$

$$\sum_{i=0}^{n} ic^{i} = \frac{nc^{n+2} - (n+1)c^{n+1} + c}{(c-1)^{2}}, \quad c \neq 1$$

$$\sum_{i=0}^{\infty} ic^{i} = \frac{c}{(1-c)^{2}}, \quad |c| < 1$$

\subsubsection{Binomial Identities}

\begin{itemize}[leftmargin=*, nosep]
    \item $\binom{n}{k} = \frac{n}{k}\binom{n-1}{k-1} = \binom{n-1}{k-1} + \binom{n-1}{k}$
    \item $\binom{n}{h}\binom{n-h}{k} = \binom{n}{k}\binom{n-k}{h}$ \hfill $\binom{n}{k} = \frac{n+1-k}{k} \binom{n}{k-1}$
    \item $\sum k \binom{n}{k} = n 2^{n-1}$ \hfill $\sum k^2 \binom{n}{k} = (n+n^2)2^{n-2}$
    \item $\sum_{j=0}^k \binom{m}{j}\binom{n-m}{k-j} = \binom{n}{k}$ \hfill $\sum \binom{m}{j}^2 = \binom{2m}{m}$
    \item $\sum_{m=k}^n \binom{m}{k} = \binom{n+1}{k+1}$ \hfill $\sum_{k=0}^{\lfloor n/2 \rfloor} \binom{n-k}{k} = F_{n+1}$
    \item $(x+y)^n = \sum \binom{n}{k}x^{n-k}y^k$ \hfill $(1+x)^n = \sum \binom{n}{k}x^k$
\end{itemize}


\subsubsection{Progressão Aritmética (PA)}

\textbf{Termo Geral:} \quad $ a_n = a_1 + (n-1)r $

\textbf{Soma dos Termos:} \quad $ S_n = \frac{(a_1 + a_n)n}{2} $

\textbf{PA de 2ª Ordem:} \quad $ a_n = An^2 + Bn + C $

\subsubsection{Progressão Geométrica (PG)}

\textbf{Termo Geral:} \quad $ a_n = a_1 \cdot q^{n-1} $

\textbf{Soma Finita:} \quad $ S_n = a_1 \frac{q^n - 1}{q - 1} $

\textbf{Soma Infinita ($|q|<1$):} \quad $ S_{\infty} = \frac{a_1}{1 - q} $


\subsubsection{Funções Geradoras}

\begin{align*}
    (1+x)^n &= \sum_{k\ge 0} \binom{n}{k} x^k \\[6pt]
    \frac{1}{(1-x)^{m+1}} &= \sum_{k\ge 0} \binom{k+m}{m} x^k \\[6pt]
    \frac{x^m}{(1-x)^{m+1}} &= \sum_{k\ge 0} \binom{k}{m} x^k
\end{align*}

\subsection{Number Theory}

\subsubsection{Identities}

$$ \sum_{d|n} \varphi(d) = n$$

$$ \sum_{\substack{i < n \\ \text{gcd}(i, n) = 1}} i = n \cdot \frac{\varphi(n)}{2} $$

$$ |\{(x, y) : \, 1 \leq x, y \leq n, \, \text{gcd}(x, y) = 1\}| = \sum_{d = 1}^n \mu(d) \left \lfloor \frac{n}{d} \right \rfloor^2$$

$$\sum_{x = 1}^n \sum_{y = 1}^n \text{gcd}(x, y)
= \sum_{k = 1}^n k \sum_{k|l}^{n} \left \lfloor \frac{n}{l} \right \rfloor^2 \mu\left(\frac{l}{k}\right)$$

$$\sum_{x = 1}^n \sum_{y = x}^n \text{gcd}(x, y) = \sum_{g = 1}^n \sum_{g|d}^n g \cdot \varphi\left(\frac{d}{g}\right)$$

$$\sum_{x = 1}^n \sum_{y = 1}^n \text{lcm}(x, y)
= \sum_{d = 1}^n d \, \mu(d) \sum_{d|l}^{n} l \binom{\lfloor \frac{n}{l} \rfloor + 1}{2}^2$$

$$\sum_{x = 1}^n \sum_{y = x+1}^n \text{lcm}(x, y) = \sum_{g = 1}^n \sum_{g|d}^n d \cdot \varphi\left(\frac{d}{g}\right) \cdot \frac{d}{g} \cdot \frac 1 2$$

$$\sum_{x \in A} \sum_{y \in A} \text{gcd}(x, y)
= \sum_{t = 1}^n \left(\sum_{l | t} \frac{t}{l} \mu(l)\right) \left(\sum_{t|a} \text{freq}[a]\right)^2$$

$$\sum_{x \in A} \sum_{y \in A} \text{lcm}(x, y)
= \sum_{t = 1}^n \left(\sum_{l | t} \frac{l}{t} \mu(l)\right) \left(\sum_{a \in A, \, t|a} a\right)^2$$

\subsubsection{Large Prime Gaps}
For numbers until $10^9$ the largest gap is 400.\\
For numbers until $10^{18}$ the largest gap is 1500.\\[0.5cm]

\subsubsection{Prime counting function - \texorpdfstring{$\pi(x)$}{}} The prime counting function is asymptotic to $\frac{x}{\log x}$, by the prime number theorem.

\

\begin{tabular}{|c|c|c|c|c|c|c|c|c|}
\hline
  \cellcolor{gray!40} x&10&$10^2$&$10^3$&$10^4$&$10^5$&$10^6$&$10^7$\\ \hline
  \cellcolor{gray!40} $\pi(x)$& 4 & 25 & 168 & 1\,229 & 9\,592 & 78\,498 & 664\,579\\ \hline
\end{tabular}

\

\subsubsection{Some Primes}

999999937 $\quad$ 1000000007 $\quad$ 1000000009 $\quad$ 1000000021 $\quad$ 1000000033
$10^{18} - 11 \quad\quad 10^{18} + 3 \quad\quad\quad 2305843009213693951 = 2^{61} - 1$

\subsubsection{Number of Divisors}

\begin{table}[H]
    \centering
    \begin{tabular}{|c|c|c|c|c|c|c|c|c|c|c|c|}
        \hline
        \cellcolor{gray!40} $n$ & 6 & 60 & 360 & 5040 & 55440 & 720720 & 4324320 \\
        \hline
        \cellcolor{gray!40} $d(n)$ & 4 & 12 & 24 & 60 & 120 & 240 & 384\\
        \hline
    \end{tabular}
\end{table}

\vspace{-30pt}

%  367567200 d(n) = 1152
%  6983776800 d(n) = 2304
%  13967553600 d(n) = 2688
%  321253732800 d(n) = 5376
%  18401055938125660800 d(n) = 184320

\begin{table}[H]
    \centering
    \begin{tabular}{|c|c|c|c|c|}
        \hline
        \cellcolor{gray!40} $n$ & 367567200 & 6983776800 & 13967553600 & 321253732800
        \\
        \hline
        \cellcolor{gray!40} $d(n)$ & 1152 & 2304 &  2688 &
        5376 \\
        \hline
    \end{tabular}
\end{table}

18401055938125660800 $\approx 2\text{e}18$ is   highly composite with 184320 divisors.

For numbers up to $10^{88}$, $d(n) < 3.6 \sqrt[3]{n}$.

\subsubsection{Lucas's Theorem}

$$ \binom{n}{m} \equiv \prod_{i=0}^k \binom{n_i}{m_i} \quad (\text{mod } p) $$

For $p$ prime. $n_i$ and $m_i$ are the coefficients of the representations of $n$    and $m$ in base $p$. In particular, $\binom{n}{m}$ is odd if and only if $n$ is a submask of $m$.

\subsubsection{Fermat's Theorems}
Let $p$ be a prime number and $a$ an integer, then:
$$a^p \equiv a \quad (\text{mod } p)$$
$$a^{p-1} \equiv 1 \quad (\text{mod } p)$$

\textbf{Lemma:} Let $p$ be a prime number and $a$ and $b$ integers, then:
$$(a+b)^{p} \equiv a^{p} + b^{p} \quad (\text{mod } p)$$

\textbf{Lemma:} Let $p$ be a prime number and $a$ an integer. The inverse of $a$ modulo $p$ is $a^{p-2}$:

$$a^{-1} \equiv a^{p-2} \quad (\text{mod } p)$$

\subsubsection{Taking modulo at the exponent}

If $a$ and $m$ are coprime, then

$$ a^n \equiv a^{n \text{ mod } \varphi(m)} \quad (\text{mod } m) $$

Generally, if $n \geq \log_2 m$, then

$$ a^n \equiv a^{\varphi(m) + [n \text{ mod } \varphi(m)]} \quad (\text{mod } m) $$

\subsubsection{Mobius invertion}

If $g(n) = \sum_{d|n} f(d)$, then $f(n) = \sum_{d|n} g(d) \mu(\frac{n}{d})$.

A more useful definition is: $\sum_{d|n} \mu(d) = [n = 1]$

Example, sum of LCM:

\begin{align*}
    \sum_{i = 1}^n \sum_{j = 1}^n \text{lcm}(i, j) &=
    \sum_{k = 1}^n\sum_{i=1}^n\sum_{j=1}^n [\text{gcd}(i, j) = k] \frac{ij}{k} \\
    &= \sum_{k = 1}^n\sum_{a=1}^{\lfloor \frac{n}{k} \rfloor}\sum_{b=1}^{\lfloor \frac{n}{k} \rfloor} [\text{gcd}(a, b) = 1] abk \\
    &= \sum_{k = 1}^n k \sum_{a=1}^{\lfloor \frac{n}{k} \rfloor} a \sum_{b=1}^{\lfloor \frac{n}{k} \rfloor} b \sum_{d=1}^{\lfloor \frac{n}{k} \rfloor}[d|a]\,[d|b] \, \mu(d) \\
    &= \sum_{k = 1}^n k \sum_{d=1}^{\lfloor \frac{n}{k} \rfloor}  \mu(d) \left(\sum_{a=1}^{\lfloor \frac{n}{k} \rfloor} [d|a] \, a\right) \left(\sum_{b=1}^{\lfloor \frac{n}{k} \rfloor} [d|b] \, b\right) \\
    &= \sum_{k = 1}^n k \sum_{d=1}^{\lfloor \frac{n}{k} \rfloor}  \mu(d) \left(\sum_{p=1}^{\lfloor \frac{n}{kd} \rfloor} p\right) \left(\sum_{q=1}^{\lfloor \frac{n}{kd} \rfloor} q\right) \\
\end{align*}

\subsubsection{Chicken McNugget Theorem}

Given two \textbf{coprime} numbers $n, m$, the largest number that cannot be written as a linear combination of them is $nm - n - m$.

\begin{itemize}
    \item There are $\frac{(n-1)(m-1)}{2}$ non-negative integers that cannot be written as a linear combination of $n$ and $m$;
    \item For each pair $(k, nm - n - m - k)$, for $k \geq 0$, exactly one can be written.
\end{itemize}


\subsubsection{Harmonic Lemma}

This technique computes sums of the form $\sum_{i=1}^{n} f\left(\left\lfloor \frac{n}{i} \right\rfloor\right)$ in $O(\sqrt{n})$.

The value of $\lfloor \frac{n}{i} \rfloor$ is constant over blocks. If we are at index $l$, the value $v = \lfloor \frac{n}{l} \rfloor$ remains constant up to index $r = \lfloor \frac{n}{v} \rfloor$. We can iterate through these blocks instead of individual indices.

\begin{verbatim}
long long sum = 0;
for (int l = 1, r; l <= n; l = r + 1) {
    int val = n / l;
    r = n / val;
    sum += (long long)(r - l + 1) * f(val);
}
\end{verbatim}

\textbf{Generalization:} For sums like $\sum_{i=1}^{n} g(i) \cdot f(\lfloor \frac{n}{i} \rfloor)$, the contribution of a block $[l, r]$ is:
$$ (G(r) - G(l-1)) \cdot f\left(\left\lfloor \frac{n}{l} \right\rfloor\right) $$
where $G(k)$ is the prefix sum $\sum_{i=1}^k g(i)$, which must be efficiently computable.

\subsubsection{Goldbach's Conjecture}
Every even integer greater than 2 is the sum of two prime numbers. Verified for all even numbers up to $4 \times 10^{18}$

\subsection{Combinatória}

\subsubsection{Stars and Bars Bounded}

De quantas formas podemos definir uma sequencia $x$ tal que:

$$ \sum^{n}_{i=0} x_i = M $$
$$ l_i \leq x_i \leq r_i $$

\paragraph{Existência de solução}

A existência  de solução é garantida se e somente se:

$$ \sum^{n}_{i=0} l_i \leq M \leq \sum^{n}_{i=0} r_i $$

\paragraph{Solução com DP em $ O(NM^2) $}

$$ DP(N, M) = \sum^{r_n}_{i=l_n} DP(N-1, M-i) $$

$$
DP(0, M) =
    \begin{cases}
        0, M \neq 0 \\
        1, M = 0
    \end{cases}
$$

\paragraph{Caso Particular $ l_i = 0 , r_i = M$}

Stars and Bars com N-1 Barras

$$ Resposta = \binom{M+N-1}{M} $$

\paragraph{Redução do problema para $l_i = 0$}

$$ M \xleftarrow{} M - \sum^{n}_{i=0} l_i $$
$$ x_i \xleftarrow{} x_i - l_i $$
$$ r_i \xleftarrow{} r_i - l_i $$
$$ l_i \xleftarrow{} 0 $$

\paragraph{Caso Especial $ r_i = R $}

$$ Resposta = \sum^{N}_{k=0} (-1)^k \binom{N}{k} \binom{M - (R + 1) k + N - 1}{N-1} $$

\subsubsection{Formulas Básicas}

\subsubsection*{Permutação Circular}
O número de maneiras de ordenar $n$ objetos distintos em um círculo.
$$ PC_n = (n-1)! $$

\subsubsection*{Permutação com Repetição}
O número de maneiras de ordenar $n$ objetos, onde existem $n_1$ objetos idênticos de um tipo, $n_2$ de outro, e assim por diante.
$$ P_n^{n_1, n_2, \dots, n_k} = \frac{n!}{n_1! n_2! \dots n_k!} $$

\subsubsection*{Combinação Caótica (Desarranjo)}
O número de permutações de $n$ objetos onde nenhum objeto permanece em sua posição original.
$$ D_n = n! \sum_{k=0}^{n} \frac{(-1)^k}{k!} $$
Aproximação para $n$ grande:
$$ D_n \approx \frac{n!}{e} $$

\subsubsection*{Generalization for calculating number of elements in exactly \( r \) sets}
A princípio da inclusão-exclusão pode ser reescrita para calcular o número de elementos que estão presentes em zero conjuntos:
\[
\left| \bigcap_{i=1}^{n} A_i \right| = \sum_{m=0}^{n} (-1)^m \sum_{\substack{|X|=m}} \left| \bigcap_{i \in X} A_i \right|
\]
Considere sua generalização para calcular o número de elementos que estão presentes em exatamente \( r \) conjuntos:
\[
\left| \bigcup_{|B|=r} \left[ \bigcap_{i \in B} A_i \cap \bigcap_{j \notin B} A_j \right] \right| = \sum_{m=r}^{n} (-1)^{m-r} \binom{m}{r} \sum_{\substack{|X|=m}} \left| \bigcap_{i \in X} A_i \right|
\]

\begin{itemize}[leftmargin=*, nosep]
    \item \textbf{Catalan:} $C_n = \frac{1}{n+1}\binom{2n}{n}$. $1, 1, 2, 5, 14, 42$.
    \item \textbf{Stirling 2:} $S(n,k) = k S(n-1, k) + S(n-1, k-1)$.
    \item \textbf{Vandermonde:} $\sum_{k=0}^r \binom{n}{k}\binom{m}{r-k} = \binom{n+m}{r}$.
    \item \textbf{Hockey-Stick:} $\sum_{i=r}^n \binom{i}{r} = \binom{n+1}{r+1}$.
\end{itemize}
